%% oak.tex -- "One huge oak": the theme=oak margin describes itself. LuaLaTeX:
%%   lualatex oak.tex
\documentclass[11pt,a4paper]{article}
\renewcommand\rmdefault{cmr}\renewcommand\sfdefault{cmss}\renewcommand\ttdefault{cmtt}
\usepackage[T1]{fontenc}
\usepackage[a4paper,left=68mm,right=22mm,top=32mm,bottom=30mm,headsep=8mm]{geometry}
\usepackage{fancyhdr}
\usepackage[grown,theme=oak]{florilegium}

\pagestyle{fancy}\fancyhf{}
\fancyhead[R]{\small\textsc{One huge oak}}
\fancyfoot[C]{\small\thepage}
\renewcommand\headrulewidth{0pt}
\newcommand\code[1]{\texttt{\small #1}}

\begin{document}
\section{A tree for every page}

\florilegiuminitial[4]{T}{his} page has an oak on it, and the letter that opens this paragraph
is an oak too, its shape filled with bark and a crown growing out of its top. The package draws
them with \code{theme=oak}. Nothing is stored: every plate of bark, every leaf and every cat is
computed while the page is written, and each page grows a different tree.

The tree does not go in the margin so much as around the text. The program is told where the
text block is and treats it as something the tree cannot grow through, as a real tree treats a
wall. The trunk stands where there is the most room, and it is only as thick as that room
allows. The limbs grow toward points scattered over the free parts of the page, each point
pulling on the nearest branch tip, so the crown fills what the words leave and stops where
they start.

\section{What lives in it}

Cats take the perches a cat would choose, birds the limbs that are nearly level, and nests the
forks. Butterflies and bees keep to the open air, clear of the tree. There may be a swing, a
heart cut into the trunk long ago, a few mushrooms in the grass.

\section{The seasons}

The command \code{\string\florilegiumage} chooses the season: 0 is spring, 0.5 summer and 1
autumn, and past the package's usual range 1.5 is winter and 2 a dead tree. The option
\code{aging} takes a document from spring to autumn, one page at a time.

\newpage
\section{Where the text stops}

On the last page the tree notices where the words end, and takes the room below them too.

\florilegiumage{1}
Set \code{\string\florilegiumtextend} after the last line, as here.
\florilegiumtextend
\end{document}
